\section{Why Reliability of AI Agents Matters for Cybersecurity}

The reliability of AI agents is not merely an academic concern: it is an
operational safety requirement. When an AI agent operates autonomously in a
SOC, its outputs directly influence incident response decisions that affect
the security posture of entire organisations. Unlike traditional software,
where bugs produce deterministic incorrect outputs, LLM-based agents
introduce a qualitatively different failure surface: \textit{stochastic,
context-dependent, and confidence-inflated errors} that are difficult to
detect and dangerous to act upon.

\textbf{Risks of autonomous AI decisions.}
A human analyst who encounters an unfamiliar threat pattern will typically
escalate or seek additional context. An LLM-based agent, by contrast, may
confidently classify a novel threat using superficially similar patterns from
its training data, a failure mode known as hallucination. In cybersecurity,
hallucination is particularly dangerous because the domain is dense with
specific factual claims (CVE identifiers, threat actor names, TTP chains)
that are easy to fabricate convincingly but difficult to verify in
real time~\cite{llmagent_halluc_survey2025}. The consequences of
incorrect autonomous decisions are often irreversible: once a compromised
host is left un-isolated or a fabricated IOC enters a shared feed, the damage
propagates faster than human review can contain it.

\textbf{Impact on incident response.}
Consider a three-stage pipeline: triage, enrichment, and remediation.
If triage misclassifies a ransomware intrusion as a benign false
positive, enrichment never investigates and remediation recommends no
action, so the intrusion persists and exfiltration proceeds undetected.
If instead triage hallucinates a threat-actor attribution, for example
pinning a commodity phishing campaign on an APT group, enrichment may
propagate fabricated IOCs to shared STIX/TAXII feeds, corrupting
defences across every organisation that consumes them. A third agent
lacking EDR containment access (capability mismatch) may identify the
threat correctly yet fail to execute isolation, introducing a critical
response delay.

\textbf{Trust and accountability.}
SOC analysts must trust agent outputs enough to act on them, but not so
much that they stop verifying. Agents that fail silently, plausible but
wrong, erode that calibration: repeated false confidence pushes analysts
into either over-reliance (automation complacency) or disengagement
(automation aversion), both degrading effectiveness. Under regulatory
frameworks such as NIS2 or DORA the challenge compounds, since
incident-response decisions must be traceable whether a human or an agent
made the call.

\textbf{Why hallucination is uniquely dangerous in SOC environments.}
Hallucination in security operations is \textit{actionable}. A hallucinated
medical diagnosis triggers further tests; a hallucinated legal citation is
caught in review; but a hallucinated threat attribution can trigger immediate
containment (host isolation, feed updates) with real cost and often no
verification checkpoint before execution. High confidence, factual density,
and automated downstream actions make SOC hallucination uniquely dangerous
among AI application domains.

\textbf{Why existing assurance approaches fall short.}
Classical software testing does not catch stochastic, context-dependent
agent failures: a test suite that passes at build time says nothing about
reliability under a new CVE or novel ransomware variant. SOC peer review
is the safety net for sensitive decisions, but autonomous agents are
deployed precisely to remove that bottleneck. Adversarial ML benchmarks
measure robustness against crafted perturbations rather than against the
benign-but-unfamiliar inputs that dominate SOC workload. Hallucination
benchmarks~\cite{bang2025hallulens} count how often models fabricate
facts but do not connect fabrication to the downstream security actions
that turn a wrong answer into operational harm. These gaps motivate a
structured framework for measuring agent reliability \textit{before
deployment} and monitoring it \textit{during operation}, which is what
\ARES{} provides.
